\section{Scope and critical reading of the corpus}
\label{sec:corpus}

\subsection{Review question and boundary}

The purpose of this review is not to assemble every result involving
HBM, CPUs, or language models.  It is to determine which claims in the
supplied corpus can support a thesis experiment on the CRESCO8 Xeon
Max~9480 nodes.  Each source is read against three questions:

\begin{enumerate}
  \item What did the work actually measure or implement?
  \item Which variables were changed together, and therefore cannot be
        separated causally?
  \item Which part of the result transfers to quantized, CPU-only LLM
        inference on CRESCO8?
\end{enumerate}

This is a \emph{corpus-bounded} review.  It can identify questions not
answered by the four supplied works, but it cannot by itself establish
that no other publication has answered them.  Any eventual novelty
claim in the thesis should therefore follow a systematic, dated search
of the wider literature and software landscape.

\subsection{Source verification}

The local file
\texttt{Lessons from HPL and HPL-MxP on Aurora.pdf} now contains the
intended paper, \emph{Sustaining Exascale Performance: Lessons from HPL
and HPL-MxP on Aurora}, by Goto et al., arXiv:2604.09517v1
~\cite{goto2026aurorahpl}.  Its title, authors, twelve-page extent and
reported campaign results agree with the paper itself.

The four works occupy different evidential roles.  Table
\ref{tab:corpus-roles} prevents results from one role being used as if
they answered another.

\begin{table}[htbp]
\centering
\small
\caption{Role and transfer limit of each reviewed work.}
\label{tab:corpus-roles}
\begin{tabularx}{\textwidth}{@{}P{2.55cm}YYP{3.35cm}@{}}
\toprule
\textbf{Work} & \textbf{Direct contribution} &
\textbf{Most useful evidence here} & \textbf{Principal transfer limit} \\
\midrule
Na et al.\ \cite{na2024llmcpu}
& Xeon Max LLM inference measurements and CPU--GPU-offload comparison
& Workload behaviour, AMX/HBM-era CPU viability, NUMA sensitivity
& BF16 IPEX experiments; not a controlled HBM--DDR comparison and not
  quantized \lcpp{} \\
\addlinespace
Ibeid et al.\ \cite{ibeid2025aurora}
& Xeon Max memory modes, microbenchmarks and three HPC applications
& Achievable HBM/DDR bandwidth, latency, capacity and locality mechanisms
& Application and mode results are not transformer measurements; much
  of the mode matrix comes from Borealis, not production Aurora \\
\addlinespace
Goto et al.\ \cite{goto2026aurorahpl}
& Deployment experience for HPL and HPL-MxP at near-full Aurora scale
& Explicit tier staging, deterministic resource mapping, phase-specific
  orchestration
& GPU-dominated heterogeneous benchmarks; impact classes are engineering
  assessments rather than controlled ablations \\
\addlinespace
Dongarra et al.\ \cite{dongarra2026gpus}
& Architectural argument for matrix-enhanced, HBM-equipped CPUs
& Motivation and hypotheses about data movement and software maturity
& Position paper with partly projected companion results; not a Max~9480
  baseline \\
\bottomrule
\end{tabularx}
\end{table}

\subsection{Na et al.: the closest workload precedent}
\label{sec:na-critical}

Na et al.\ evaluate BF16 inference with Intel Extension for PyTorch on
a dual-socket Xeon Max~9468 system and compare it with an Ice Lake
server and single A100 and H100 GPUs running FlexGen
~\cite{na2024llmcpu}.  Their main Xeon Max experiments bind one
48-core socket; a separate core-count experiment extends to 96 cores.
The study covers OPT and LLaMA-2 models, prompt length 128, generation
length 32, and several batch sizes, with an additional prompt-length
sweep.

Three achievements are directly relevant.

\begin{itemize}
  \item Xeon Max improves geometric-mean throughput by 3.2--6.3 times
        over the tested Ice Lake system, depending on batch size.
        This establishes that the combined change in matrix engine,
        memory system, cores, microarchitecture and software can be
        substantial.  It does not isolate the contribution of HBM or
        AMX.
  \item Quadrant+Flat is the best average memory configuration in the
        reported setup, while unmanaged SNC4 performs poorly.  This
        makes placement a first-order experimental variable, but the
        paper does not disclose enough first-touch and page-residency
        detail to reproduce the causal mechanism exactly.
  \item A single 48-core socket is faster on average than the naive
        96-core execution.  The result demonstrates a locality failure;
        it does not show that two sockets cannot help when used as
        explicitly bound replicas or with topology-aware partitioning.
\end{itemize}

The CPU--GPU comparison also needs a narrow reading.  At low batch and
for models that require host offload, the direct CPU execution avoids
repeated PCIe movement and can outperform the tested FlexGen path.  The
result reverses in at least one larger-batch, longer-prompt case:
for LLaMA2-70B at batch 16 and prompt length at least 256, the H100
offload run reports lower latency.  Models that fit wholly in GPU memory
also favour the GPUs.  The paper therefore supports a capacity and
data-movement argument, not a general ranking of CPUs and GPUs.

For the present thesis, the largest omissions are quantized weights,
long sustained decode, explicit separation of time to first token and
inter-token latency, verified memory residency, uncertainty, energy and
online serving behaviour.  These omissions define useful measurements,
but not all of them need to become thesis contributions.

\subsection{Ibeid et al.: the mechanism of the memory system}
\label{sec:ibeid-critical}

Ibeid et al.\ characterise Xeon Max memory through STREAM, Intel Memory
Latency Checker, PCIe and MPI tests, and three HPC applications
~\cite{ibeid2025aurora}.  Their Borealis testbed exposes the full
combination of Flat/Cache and Quadrant/SNC4 modes.  Production Aurora
was in Quadrant+Flat for the reported application study, so the two
sources of evidence must not be conflated.

The clearest quantitative result is the single-socket STREAM Triad
comparison:

\begin{table}[htbp]
\centering
\small
\caption{Borealis single-socket STREAM Triad bandwidth, in GB/s
~\cite{ibeid2025aurora}.}
\label{tab:stream-critical}
\begin{tabular}{@{}lccc@{}}
\toprule
& \multicolumn{2}{c}{\textbf{Flat mode}} & \textbf{Cache mode} \\
\cmidrule(lr){2-3}
\textbf{Clustering} & \textbf{DDR5} & \textbf{HBM} & \textbf{DDR5+HBM cache} \\
\midrule
Quadrant, full socket & 235 & 680 & 560 \\
SNC4, full socket     & 245 & 840 & 575 \\
SNC4, one quadrant    &  60 & 210 & 150 \\
\bottomrule
\end{tabular}
\end{table}

For the two full-socket Flat rows, HBM sustains approximately
2.9--3.4 times the DDR5 bandwidth.  SNC4 raises full-socket HBM
bandwidth from 680 to 840\,GB/s when the benchmark is deliberately
placed.  At low load, however, the paper reports lower unloaded latency
for Quadrant/DDR5 (112.9\,ns) than for Quadrant/HBM (134.9\,ns).  The
useful conclusion is conditional: HBM is attractive for a concurrent,
bandwidth-intensive live set that fits; DDR5 provides capacity and can
have a latency advantage under light load; locality determines how
much of either resource is obtained.

The application results reinforce the conditional nature of memory
mode.  Explicit HBM is strongest when the working set fits, Cache mode
can help beyond HBM capacity, and direct-mapped cache conflicts can
distort the boundary.  The authors use page shuffling (``Fake NUMA'')
to reduce those conflicts.  Consequently, a result close to 64\,GB
cannot be interpreted from nominal capacity alone.

This paper supplies a mechanism and a validation target for CRESCO8:
measure the local HBM and DDR bandwidth actually available under the
same binding used for inference.  Its STREAM ratio is not an expected
LLM speedup.  The inference runtime has additional computation, cache
reuse, synchronization and non-streaming accesses.

\subsection{Goto et al.: production tiering and orchestration}
\label{sec:goto-critical}

Goto et al.\ describe three Aurora deployment campaigns
~\cite{goto2026aurorahpl}.  FP64 HPL progressed from 0.585\,EF/s on
5,439 nodes at SC23 to 1.012\,EF/s on 9,234 nodes, with 78.8\%
parallel scaling efficiency.  HPL-MxP reached 10.6\,EF/s on 9,234
nodes at ISC24 and 11.64\,EF/s on 9,500 nodes at SC24.  The reported
11.5-fold relation between HPL-MxP and HPL combines a change in
precision, algorithm, process grid and node count; it is not a
like-for-like processor speedup.

The corrected paper adds a particularly relevant memory-management
example.  Each Aurora CPU has 64\,GB HBM2e and 512\,GB DDR5.  Because
the HPL global matrix exceeds HBM capacity, CPU-managed bulk storage
remains in DDR.  Tiles needed for panel factorisation, TRSM and row
exchange are staged into HBM.  The design is therefore an observed,
production-scale instance of \emph{capacity in DDR plus selected
bandwidth-critical data in HBM}.  Overall runtime remains dominated by
GPU GEMM, so the paper neither quantifies the isolated benefit of this
staging nor predicts an LLM result.

The paper's second transferable result is methodological.  Processes
are co-located with CPU cores, GPU, memory and NIC on the same CPU/NUMA
domain; CPU--GPU transfers are overlapped explicitly; and NICs are
assigned by communication phase.  A hybrid point-to-point/collective
path improves resilience to transient network events.  These details
show that resource mapping and run-time configuration are part of the
experimental treatment, not incidental setup.

Table~II in the paper labels NUMA mapping, CPU--GPU overlap and AMX as
``critical'', among other classifications.  Its footnote states that
these labels are engineering assessments rather than controlled
ablations.  The SC24 HPL-MxP result is about 9.8\% above the ISC24
result, and the authors identify AMX enablement as the primary
configuration change.  This campaign comparison is valuable operational
evidence, but it does not exclude all contemporaneous tuning changes.
The safe conclusion is that AMX contributed to the production result,
not that the paper measures a portable 9.8\% AMX effect.

\subsection{Dongarra et al.: a research position}
\label{sec:dongarra-critical}

Dongarra, Hoefler and Matsuoka argue that CPUs combining matrix engines,
low-precision arithmetic and on-package HBM reduce the arithmetic,
bandwidth and data-movement disadvantages that motivated discrete GPUs
~\cite{dongarra2026gpus}.  They retain an important boundary: GPUs
remain central to frontier training, while HBM-equipped CPUs may be
attractive for inference and for workflows that mix dense and irregular
operations.

The manuscript is useful for forming hypotheses, especially that
software maturity and explicit memory placement may limit performance.
It is not a direct baseline for CRESCO8.  Its empirical argument relies
on a companion study described as forthcoming: A64FX decode results are
measured, while parts of the newer matrix-enhanced CPU discussion,
including prefill and energy, are projected from specifications.  The
workload scale and formats also differ markedly from a practical
single-node Max~9480 study.

\subsection{Assessment}

The corpus is coherent once the roles are separated.  Na et al.\ provide
the nearest workload evidence; Ibeid et al.\ explain the memory
mechanisms; Goto et al.\ demonstrate deliberate tiering and end-to-end
resource mapping at production scale; and Dongarra et al.\ motivate why
the problem matters.  None of the four performs the same-socket,
same-runtime HBM-versus-DDR experiment required to attribute a
quantized LLM result to memory tier.
